\chapter{Fondamenti Metodologici}
\label{cap2}
% Machine learning / Deep learning / 
\section{Machine Learning}
Il \textbf{machine learning} è un approccio alla risoluzione di problemi di riconoscimento di pattern in cui, anziché codificare manualmente regole esplicite, si utilizza un ampio insieme di dati di addestramento per determinare i parametri di un modello adattivo. Sostanzialmente il \textit{machine learning} impiega la costruzione di modelli matematici necessari ad una maggiore comprensione dei dati. Tali modelli vengono programmati, attraverso dei parametri, ed esposti ai dati, per poi poter essere usati nella comprensione di nuove informazioni \cite{vanderplas2016python}. 
Esistono due categorie principali di \textit{machine learning}: apprendimento supervisionato e non supervisionato.
L'\textbf{apprendimento supervisionato} richiede di modellare la relazione tra le \textit{feature} dei dati e le etichette associate, con l'obiettivo di generalizzare e sviluppare la capacità di classificare correttamente esempi diversi da quelli usati in fase di addestramento.
Si articola in due sottocategorie principali in funzione della natura dell'output da predire \cite{bishop2006pattern}. Nella classificazione (\textit{classification}), le etichette sono categorie discrete, ovvero un insieme finito di classi predefinite, a cui il modello deve assegnare ogni esempio di input. Un esempio tipico è il riconoscimento di azioni in un video, dove l'output è una tra le etichette disponibili (es. \textit{``take''}, \textit{``cut''}, \textit{``pour''}). Nella regressione (\textit{regression}), l'output è invece una variabile continua, ovvero il modello deve stimare un valore numerico, come la durata stimata di un'azione.
Il processo di apprendimento consiste, inoltre, nell'aggiustare iterativamente i parametri del modello fino a minimizzare la \textit{loss function}, la quale misura quanto le predizioni del modello si discostano dalle etichette corrette \cite{bishop2006pattern}.
L'\textbf{apprendimento non supervisionato}, al contrario, consiste nel modellare la struttura di un dataset in assenza di etichette di riferimento. L'obiettivo può essere individuare gruppi di esempi simili all'interno dei dati (\textit{clustering}), determinare la distribuzione dei dati nello spazio di input (\textit{density estimation}), o ridurne la dimensionalità preservando le informazioni rilevanti \cite{bishop2006pattern}.
Una terza categoria, distinta dalle precedenti, è l'\textbf{apprendimento per rinforzo} (\textit{reinforcement learning}): anziché apprendere da un insieme fisso di dati etichettati o non etichettati, il modello impara interagendo con un ambiente, ricevendo un segnale di ricompensa (\textit{reward}) in risposta alle proprie azioni e adattando il proprio comportamento per massimizzarlo nel tempo. \cite{bishop2006pattern, sutton2018}. A differenza dell'apprendimento supervisionato, non esiste una risposta corretta esplicita; il modello deve scoprire autonomamente quale sequenza di azioni conduce al risultato desiderato.

\section{Deep Learning}
\subsection{Reti neurali e deep learning}
Il \textbf{deep learning} è una branca del machine learning basata sull'utilizzo di reti neurali artificiali con molteplici strati di elaborazione \cite{goodfellow2016deep}. Una \textbf{rete neurale artificiale} è un modello  computazionale ispirato al funzionamento del cervello 
biologico, composto da unità elementari chiamate \textbf{neuroni artificiali} organizzate in strati 
\cite{goodfellow2016deep}. Ogni strato apprende rappresentazioni progressivamente più astratte dell'input: i primi strati catturano caratteristiche elementari (bordi, texture), mentre gli strati più profondi combinano queste informazioni in rappresentazioni semantiche significative \cite{goodfellow2016deep}. Questo processo di \textit{rappresentation learning} ha reso il \textit{deep learning} particolarmente efficace in domini di alta dimensionalità come immagini, audio e video, dove la progettazione manuale di \textit{feature} risulterebbe impraticabile.

\subsection{Reti Convoluzionali (CNN)}
Le \textbf{reti neurali convoluzionali} (\textit{Convolutional Neural Networks}, CNN) sono architetture progettate per elaborare dati con struttura a griglia, come le immagini \cite{goodfellow2016deep}. Il loro componente chiave è il \textit{layer convoluzionale}, che applica un insieme di filtri appresi localmente sull'input, condividendo i pesi tra posizioni diverse: questo riduce drasticamente il numero di parametri rispetto a un \textit{layer} completamente connesso e introduce la capacità di riconoscere un pattern indipendentemente dalla sua posizione nell'immagine \cite{goodfellow2016deep}.
\subsection{Reti Ricorrenti (RNN e LSTM)}
Le \textbf{reti neurali ricorrenti} (\textit{Recurrent Neural Networks}, RNN) sono progettata per elaborare sequenze di dati mantenendo uno stato nascosto (\textit{hidden state}) che viene aggiornato ad ogni passo temporale, permettendo al modello di conservare memoria delle elaborazioni precedenti \cite{goodfellow2016deep}. Il loro limite principale è la difficoltà di modellare dipendenze a lungo termine, dovuta al problema del \textit{vanishing gradient}: i gradienti si attenuano propagandosi attraverso molti  passi temporali, rendendo difficile l'apprendimento di relazioni tra eventi distanti nella sequenza \cite{goodfellow2016deep}. Le \textbf{Long Short-Term Memory} (LSTM) risolvono questo problema introducendo un meccanismo di \textit{gate} che regola esplicitamente quali informazioni conservare, aggiornare o dimenticare, permettendo di modellare dipendenze temporali su orizzonti significativamente più lunghi \cite{goodfellow2016deep}. Per questo motivo le LSTM sono state ampiamente adottate nei primi modelli di action anticipation, come RU-LSTM \cite{furnari2019would, furnari2020rolling}. 
\subsection{Transformer e meccanismo di attenzione}
Lo stesso problema viene risolto dall'introduzione dei Transformer \cite{vaswani2017attention} con un approccio completamente diverso: anziché accumulare il contesto in uno stato nascosto, ogni elemento della sequenza può confrontarsi direttamente con tutti gli altri, calcolando quanto ciascuno sia rilevante per la sua interpretazione \cite{vaswani2017attention}. Questo meccanismo, detto \textit{self-attention}, permette al modello di costruire rappresentazioni contestuali dinamiche dove, per anticipare l'azione successiva, esso può considerare direttamente un'azione avvenuta tre minuti prima senza dover propagare quell'informazione attraverso tutti i passi intermedi \cite{vaswani2017attention}.
Un altro vantaggio altrettanto rilevante introdotto dai Transformer consiste nella \textbf{parallelizzazione}: mentre le RNN devono necessariamente elaborare la sequenza in ordine, il Transformer processa tutti gli elementi simultaneamente, rendendo l'addestramento significativamente più efficiente \cite{vaswani2017attention}. Questi vantaggi hanno reso i Transformer l'architettura dominante nell'\textit{action anticipation}, da modelli come AVT \cite{girdhar2021anticipative} a RAFTformer \cite{girase2023latency}, sostituendo le architetture ricorrenti nei modelli più performanti \cite{lai2024human}.
\section{Comprensione video}
\subsection{Evoluzione degli approcci}
I primi approcci al riconoscimento e all'anticipazione di azioni umane si basavano su \textit{feature handcrafted}, ovvero rappresentazioni visive progettate manualmente da esperti del dominio \cite{hu2022online}. Tra le più diffuse figurano l'\textbf{optical flow}, che codifica il movimento apparente dei pixel tra \textit{frame} consecutivi catturando informazioni dinamiche sulla traiettoria del soggetto, e descrittori basati su \textbf{dati scheletrici} (\textit{skeleton}), che rappresentano il corpo umano come un grafo e ne tracciano la posizione nel tempo \cite{hu2022online}. Questi metodi presentano limitazioni significative: le \textit{feature handcrafted} richiedono una progettazione specifica per ogni scenario applicativo, risultano sensibili a variazioni di illuminazione e a cambiamenti di punto di vista, e non generalizzano facilmente a contesti diversi da quelli per cui sono stati progettati \cite{hu2022online, lai2024human}.
L'introduzione delle reti neurali profonde ha trasformato radicalmente il campo, sostituendo le \textit{feature handcrafted} con rappresentazioni apprese direttamente dai dati, più espressive e generalizzabili. Le architetture CNN, RNN e Transformer si sono progressivamente affermate come standard per il riconoscimento e l'anticipazione di azioni \cite{lai2024human}.
\subsection{Modalità di input}
I modelli moderni per l'action anticipation non si limitano all'elaborazione del segnale RGB grezzo, ma integrano modalità di input eterogenee che arricchiscono la rappresentazione della scena \cite{lai2024human}. Il canale \textbf{RGB} fornisce informazioni sull'apparenza visiva, spesso affiancato dall'\textbf{optical flow} per catturare le dinamiche di movimento. Gli \textbf{object detections}, ovvero \textit{bounding box} e categorie degli oggetti rilevati da \textit{detector} pre-addestrati, aggiungono informazione semantica sulla stato della scena \cite{lai2024human}. I \textbf{past actions segments}, ovvero le etichette testuali delle azioni già eseguite, sono utili come contesto storico, specialmente su dataset a struttura sequenziale regolare come Breakfast Actions e 50 Salads \cite{sener2020temporal}. Infine, \textbf{hand masks} e \textbf{gaze} sono modalità specifiche per video egocentrici, che indicano rispettivamente la posizione delle mani e la direzione dello sguardo. La scelta delle modalità di input incide direttamente sulla confrontabilità dei modelli, poiché un sistema che sfrutta \textit{object detections} o \textit{past action segments} dispone di informazioni aggiuntive rispetto a uno che opera solo su RGB \cite{lai2024human}.
\subsection{Feature e backbone}
Nel contesto del video undestanding, il termine \textbf{feature} indica una rappresentazioni vettoriale compatta che codifica le informazioni visive rilevanti di un \textit{frame} o di un segmento video. Il \textbf{backbone} è la rete neurale responsabile dell'estrazione di queste rappresentazioni; tipicamente si tratta di architetture pre-addestrate come \textbf{I3D} \cite{carreira2017quo}, che estende le convoluzioni 2D delle CNN classiche alla dimensione temporale per catturare il movimento, oppure \textbf{ViT} (\textit{Vision Transformer}) \cite{dosovitskiy2020image}, che applica il meccanismo di \textit{self-attention} direttamente alle \textit{patch} di immagine. La scelta del \textit{backbone} ha un impatto determinante sulle performance, in quanto due modelli con architetture diverse ma con stesso \textit{backbone} possono produrre risultati più simili tra loro di quanto non facciano due modelli con la stessa architettura ma \textit{backbone} differenti \cite{lai2024human}.
\subsection{Pre-training}
Addestrare un modello di deep learning da zero richiede quantità di dati e risorse computazionali spesso proibitive. Il \textbf{pre-training} consiste nell'addestrare inizialmente il modello su un dataset di grandi dimensioni, per poi adattarlo al task specifico tramite \textit{fine-tuning} \cite{goodfellow2016deep}. Questo approccio, noto come \textit{transfer learning} permette al modello di trasferire le rappresentazioni apparse su larga scala al dominio target, ottenendo performance superiori anche con dataset di dimensioni contenute. 
\subsection{Large Language Models per la comprensione video}
I \textbf{Large Language Models} (LLM) sono modelli neurali addestrati su grandi quantità di testo, che sviluppano una rappresentazione implicita della struttura sequenziale e causale delle attività umane \cite{mittal2024can}. Applicati al task dell'\textit{action anticipation} permettono di sfruttare questa conoscenza procedurale, ad esempio che \textit{``tagliare le verdure''} segue tipicamente \textit{``lavare le verdure''}, per guidare la predizione delle azioni future, integrando il segnale visivo con una struttura semantica appresa su scala molto più ampia.